\documentclass[12pt,a4paper]{report}

% ── Codificación y lenguaje ────────────────────────────────────
\usepackage[utf8]{inputenc}
\usepackage[T1]{fontenc}
\usepackage[spanish,es-tabla]{babel}

% ── Tipografía ────────────────────────────────────────────────
\usepackage{lmodern}
\usepackage{microtype}

% ── Geometría ─────────────────────────────────────────────────
\usepackage[top=2.5cm,bottom=2.5cm,left=3cm,right=2.5cm]{geometry}

% ── Gráficos y colores ────────────────────────────────────────
\usepackage{xcolor}
\usepackage{graphicx}
\usepackage{tikz}

% ── Tablas ────────────────────────────────────────────────────
\usepackage{booktabs}
\usepackage{longtable}
\usepackage{array}
\usepackage{multirow}

% ── Código fuente ─────────────────────────────────────────────
\usepackage{listings}
\lstset{
  basicstyle=\ttfamily\small,
  keywordstyle=\color{blue},
  commentstyle=\color{gray},
  stringstyle=\color{orange!70!black},
  breaklines=true,
  frame=single,
  rulecolor=\color{black!20},
  backgroundcolor=\color{black!3},
  numbers=left,
  numberstyle=\tiny\color{gray},
  tabsize=4,
  showstringspaces=false
}

% ── Hipervínculos ─────────────────────────────────────────────
\usepackage[hidelinks,colorlinks=true,linkcolor=blue!70!black,urlcolor=blue!70!black]{hyperref}

% ── Encabezados y pies ────────────────────────────────────────
\usepackage{fancyhdr}
\pagestyle{fancy}
\fancyhf{}
\fancyhead[L]{\small HMI Clasificador de Paquetes}
\fancyhead[R]{\small \leftmark}
\fancyfoot[C]{\thepage}
\renewcommand{\headrulewidth}{0.4pt}

% ── Colores corporativos ──────────────────────────────────────
\definecolor{unerblue}{RGB}{0,75,135}
\definecolor{lightgray}{RGB}{245,245,245}
\definecolor{darkgray}{RGB}{60,60,60}

% ── Comandos utilitarios ──────────────────────────────────────
\newcommand{\cmd}[1]{\texttt{#1}}
\newcommand{\file}[1]{\texttt{#1}}
\newcommand{\class}[1]{\texttt{\textbf{#1}}}

% ─────────────────────────────────────────────────────────────
\begin{document}

% ══════════════════════════════════════════════════════════════
%  PORTADA
% ══════════════════════════════════════════════════════════════
\begin{titlepage}
  \centering
  \vspace*{1.5cm}

  {\color{unerblue}\rule{\linewidth}{3pt}}
  \vspace{0.8cm}

  {\Huge\bfseries\color{unerblue} HMI Clasificador de Paquetes}\\[0.6cm]
  {\Large Informe Técnico -- Manual de Usuario}\\[0.3cm]
  {\large Versión 1.1}

  \vspace{0.8cm}
  {\color{unerblue}\rule{\linewidth}{1pt}}

  \vspace{1.5cm}

  \begin{tabular}{ll}
    \textbf{Proyecto:}      & HMI\_Clasificador \\[4pt]
    \textbf{Institución:}   & Universidad Nacional de Entre Ríos (UNER) \\[4pt]
    \textbf{Materia:}       & Microcontroladores \\[4pt]
    \textbf{Plataforma:}    & Qt 6 / C++ -- Windows \\[4pt]
    \textbf{Repositorio:}   & \href{https://github.com/XradX314/HMI_Clasificador}{github.com/XradX314/HMI\_Clasificador} \\[4pt]
    \textbf{Fecha:}         & Mayo 2026 \\
  \end{tabular}

  \vfill
  {\color{unerblue}\rule{\linewidth}{3pt}}
\end{titlepage}

% ══════════════════════════════════════════════════════════════
%  TABLA DE CONTENIDOS
% ══════════════════════════════════════════════════════════════
\tableofcontents
\clearpage

% ══════════════════════════════════════════════════════════════
%  PARTE I – MANUAL DE USUARIO
% ══════════════════════════════════════════════════════════════
\part{Manual de Usuario}

% ──────────────────────────────────────────────────────────────
\chapter{Introducción}
% ──────────────────────────────────────────────────────────────

\section{Descripción general}

El \textbf{HMI Clasificador de Paquetes} es una aplicación de escritorio
desarrollada en Qt~6 / C++ que permite controlar y monitorear un sistema
automático de clasificación de cajas por cinta transportadora.

La comunicación con el microcontrolador (MCU) se realiza a través de un
puerto serie (UART) utilizando el protocolo binario propio \textbf{UNER}.
La aplicación ofrece:

\begin{itemize}
  \item Control de arranque, parada y reinicio del clasificador.
  \item Configuración de las salidas de clasificación (tipo de caja asignado a
        cada rama de la cinta).
  \item Modo ciego: clasificación basada en distancias configurables en lugar de
        medición activa.
  \item Calibración de alturas de referencia, tiempos de brazos y parámetros
        de hardware.
  \item Visualización en tiempo real de sensores IR, estado de los brazos
        actuadores y contadores de cajas.
  \item Indicador de heartbeat (LED animado) con monitoreo de la conexión.
  \item Temas visual claro y oscuro.
\end{itemize}

\section{Requisitos del sistema}

\begin{tabular}{ll}
  \toprule
  \textbf{Componente} & \textbf{Requisito mínimo} \\
  \midrule
  Sistema operativo & Windows 10 / 11 (64-bit) \\
  Qt Runtime & Qt 6.x (incluido en el instalador) \\
  Puerto USB & Adaptador USB-UART compatible con el MCU \\
  MCU firmware & Compatible con protocolo UNER v1.1 \\
  \bottomrule
\end{tabular}

% ──────────────────────────────────────────────────────────────
\chapter{Instalación y arranque}
% ──────────────────────────────────────────────────────────────

\section{Compilación desde código fuente}

\begin{enumerate}
  \item Instalar Qt~6 con soporte para \texttt{QSerialPort} y MinGW.
  \item Abrir \file{HMI\_Clasificador.pro} en Qt Creator.
  \item Seleccionar el kit de compilación MinGW 64-bit.
  \item Presionar \textbf{Ejecutar} (\texttt{Ctrl+R}).
\end{enumerate}

O desde línea de comandos (directorio del proyecto):

\begin{lstlisting}[language=bash,caption={Compilación por consola}]
qmake HMI_Clasificador.pro -spec win32-g++
mingw32-make -j4
\end{lstlisting}

\section{Conexión del hardware}

\begin{enumerate}
  \item Conectar el MCU al PC mediante el adaptador USB-UART.
  \item Verificar que el puerto COM aparece en el Administrador de dispositivos.
  \item Abrir la aplicación \texttt{HMI\_Clasificador.exe}.
\end{enumerate}

% ──────────────────────────────────────────────────────────────
\chapter{Interfaz gráfica}
% ──────────────────────────────────────────────────────────────

\section{Barra de herramientas}

La barra superior contiene los controles de conexión serie:

\begin{itemize}
  \item \textbf{Puerto COM:} selector del puerto serie disponible.
  \item \textbf{Baudios:} velocidad de transmisión (115200 recomendado).
  \item \textbf{Conectar / Desconectar:} establece o cierra la conexión.
  \item \textbf{Actualizar:} refresca la lista de puertos detectados.
  \item \textbf{$\circ$ Calibración:} abre el diálogo de calibración (CMD~0x63).
  \item \textbf{Ancho de caja:} abre el diálogo de ancho de referencia (CMD~0x62).
  \item \textbf{Config.\ avanzada:} abre el diálogo de parámetros de hardware (CMDs~0x64–0x67).
  \item \textbf{Modo oscuro:} alterna el tema visual.
\end{itemize}

\section{Panel Monitor}

\subsection{Heartbeat / Conexión}

Muestra el LED animado \class{LedAliveWidget} con los estados:
\begin{itemize}
  \item \textcolor{green!60!black}{\textbf{Verde}}: heartbeat recibido (online).
  \item \textbf{Gris}: sin actividad (offline).
  \item \textcolor{orange}{\textbf{Naranja}}: timeout ($> 7$ s sin heartbeat).
\end{itemize}

\subsection{Panel de control}

\begin{tabular}{lll}
  \toprule
  \textbf{Control} & \textbf{CMD} & \textbf{Descripción} \\
  \midrule
  Start & 0x50 & Inicia la cinta con los tipos de caja configurados \\
  Stop  & 0x51 & Detiene la cinta \\
  Reset & 0x53 & Reinicia el MCU y los contadores \\
  Velocidad & 0x54 & Índice de velocidad (1–10; enviado al cambiar el SpinBox) \\
  Modo Ciego & 0x60 & Activa/desactiva modo ciego + envía distancias S0→Sal0/1/2 \\
  Trigger & 0x61 & Dispara una medición manual inmediata \\
  \bottomrule
\end{tabular}

\subsection{Panel de salidas}

Tres \textit{ComboBox} permiten asignar a cada salida el tipo de caja
(Ninguna / Pequeña / Mediana / Grande).  El botón \textbf{Aplicar} envía
el CMD~0x50 con la configuración actual.

\subsection{Panel de sensores y brazos}

Muestra en tiempo real el estado de los 4 sensores IR y los 3 brazos
actuadores.  Los indicadores cambian de color al recibir los CMDs~0x5E
(IR) y~0x52 (Brazo).

\subsection{Contadores}

Muestra la cantidad de cajas entradas, salidas y en tránsito durante la
sesión actual.

\subsection{Última caja}

Muestra la altura medida (cm), el tipo inferido y el timestamp de la
última caja detectada (CMD~0x5F).

% ──────────────────────────────────────────────────────────────
\chapter{Diálogos de configuración}
% ──────────────────────────────────────────────────────────────

\section{Calibración (CMD 0x63)}

Permite ajustar las alturas de referencia del sensor HC-SR04:

\begin{itemize}
  \item \textbf{Distancia piso:} distancia del sensor al piso sin caja (\texttt{calibracion[0]}).
  \item \textbf{Caja pequeña / mediana / grande:} alturas de cada tipo (\texttt{calibracion[1..3]}).
  \item \textbf{Tolerancia:} margen de error en cm.
  \item \textbf{Extensión / Retracción:} tiempos de actuación de los brazos en ticks (1 tick = 2 ms).
\end{itemize}

Presionar \textbf{Enviar 0x63} transmite los 7 bytes al MCU.

\section{Ancho de caja (CMD 0x62)}

Permite enviar el ancho de la caja de referencia en cm.  El MCU utiliza
este valor para calcular internamente la velocidad de la cinta.

\section{Configuración avanzada (CMDs 0x64–0x67)}

Diálogo desplazable con cuatro secciones independientes:

\begin{itemize}
  \item \textbf{HC-SR04 (0x64):} pulso de trigger (µs) y timeout de eco (µs).
  \item \textbf{Servo SG90 (0x65):} período, pulso mínimo, máximo y neutro (µs).
  \item \textbf{Debounce IR (0x66):} número de ticks de debounce.
  \item \textbf{Timers / Cajas (0x67):} retractTimer, ángulos SG90, hcsrTimer,
        aliveTimer y configCajas[0..2].
\end{itemize}

Cada sección tiene su propio botón \textbf{Enviar} para transmitir
únicamente los parámetros de esa sección.

% ══════════════════════════════════════════════════════════════
%  PARTE II – DOCUMENTACIÓN TÉCNICA
% ══════════════════════════════════════════════════════════════
\part{Documentación Técnica}

% ──────────────────────────────────────────────────────────────
\chapter{Arquitectura del sistema}
% ──────────────────────────────────────────────────────────────

\section{Vista general}

El sistema se organiza en tres capas:

\begin{enumerate}
  \item \textbf{Capa de protocolo} (\file{comunicacion/}): codificación/decodificación de tramas UNER.
  \item \textbf{Capa de comunicación} (\file{comunicacion/}): gestión del puerto serie y despacho de eventos.
  \item \textbf{Capa de presentación} (\file{MainWindow.cpp}, \file{widgets/}): interfaz gráfica y lógica de usuario.
\end{enumerate}

\section{Diagrama de clases simplificado}

\begin{center}
\begin{tabular}{|l|}
\hline
\textbf{MainWindow} \\
\hline
-- SerialManager* \\
-- ConfigDialog* \\
-- VelocidadDialog* \\
-- AvanzadoDialog* \\
-- LedAliveWidget* \\
\hline
\end{tabular}
\quad $\rightarrow$ \quad
\begin{tabular}{|l|}
\hline
\textbf{SerialManager} \\
\hline
-- QSerialPort* \\
-- UnerProtocol* \\
-- QTimer* (heartbeat) \\
\hline
\end{tabular}
\quad $\rightarrow$ \quad
\begin{tabular}{|l|}
\hline
\textbf{UnerProtocol} \\
\hline
Encoder estático \\
Decoder de máquina \\
de estados \\
\hline
\end{tabular}
\end{center}

\section{Flujo de datos}

\begin{enumerate}
  \item El MCU envía bytes por UART.
  \item \class{QSerialPort} emite \texttt{readyRead} $\rightarrow$ \class{SerialManager::onDataReady}.
  \item Los bytes se pasan a \class{UnerProtocol::feed()} (máquina de estados).
  \item Al completarse una trama válida, \class{UnerProtocol} emite \texttt{frameReceived}.
  \item \class{SerialManager::dispatchFrame()} interpreta el comando y emite señales específicas
        (\texttt{cajaMedida}, \texttt{sensorIrActualizado}, \texttt{brazoActuado}, \texttt{aliveReceived}).
  \item \class{MainWindow} recibe las señales y actualiza la UI.
\end{enumerate}

% ──────────────────────────────────────────────────────────────
\chapter{Protocolo UNER}
% ──────────────────────────────────────────────────────────────

\section{Estructura de trama}

Toda trama sigue el formato:

\begin{center}
\begin{tabular}{|c|c|c|c|c|c|}
\hline
\texttt{U} & \texttt{N} & \texttt{E} & \texttt{R} & \texttt{len} & \texttt{':'} \\
\hline
\texttt{cmd} & \texttt{payload[0]} & \texttt{...} & \texttt{payload[n-1]} & \texttt{XOR} & \\
\hline
\end{tabular}
\end{center}

\begin{itemize}
  \item \textbf{Header:} 4 bytes fijos \texttt{'U','N','E','R'}.
  \item \textbf{len:} \texttt{1 + n + 1} (cmd + payload + checksum).
  \item \textbf{':'} (0x3A): token separador.
  \item \textbf{cmd:} byte de comando.
  \item \textbf{payload:} 0 a \textit{n} bytes de datos.
  \item \textbf{XOR:} XOR de todos los bytes anteriores de la trama.
\end{itemize}

\section{Tabla de comandos}

\begin{longtable}{cllp{6cm}}
  \toprule
  \textbf{CMD} & \textbf{Nombre} & \textbf{Dirección} & \textbf{Payload} \\
  \midrule
  \endhead
  0xF0 & CMD\_ALIVE & MCU$\to$HMI / HMI$\to$MCU & [0x06] (ACK bidireccional) \\
  \midrule
  0x50 & CMD\_START & HMI$\to$MCU & [0x00, tipoSal0, tipoSal1, tipoSal2] \\
  0x51 & CMD\_STOP  & HMI$\to$MCU & [0x06] \\
  0x52 & CMD\_BRAZO & MCU$\to$HMI & [máscara\_servo, estado] \\
  0x53 & CMD\_RESET & HMI$\to$MCU & [0x06] \\
  0x54 & CMD\_VELOCIDAD & HMI$\to$MCU & [velIdx × 10] (0–100) \\
  \midrule
  0x5E & CMD\_IR\_STATE  & MCU$\to$HMI & pares [outNum, IRState] \\
  0x5F & CMD\_CAJA\_DETECT & MCU$\to$HMI & [alturaCm] \\
  \midrule
  0x60 & CMD\_BLIND\_DIST & HMI$\to$MCU & [modo\_ciego, dist\_s0[0], [1], [2]] (1–4 bytes) \\
  0x61 & CMD\_TRIGGER & HMI$\to$MCU & (sin payload) \\
  0x62 & CMD\_ANCHO\_CAJA & HMI$\to$MCU & [anchoCm] \\
  0x63 & CMD\_CALIBRACION & HMI$\to$MCU & [cal0, cal1, cal2, cal3, tol, ext, ret] \\
  0x64 & CMD\_HCSR04\_CFG & HMI$\to$MCU & [trig\_pulse\_us, timeout\_hi, timeout\_lo] \\
  0x65 & CMD\_SG90\_CFG & HMI$\to$MCU & [period\_hi, lo, min\_hi, lo, max\_hi, lo, neutral\_hi, lo] \\
  0x66 & CMD\_IR\_DEBOUNCE & HMI$\to$MCU & [debounce\_ticks] \\
  0x67 & CMD\_TIMERS\_CFG & HMI$\to$MCU & [retract, angleDet, angleRep, hcsr, alive, cajas0, cajas1, cajas2] \\
  \bottomrule
\end{longtable}

\section{Verificación de checksum}

El checksum es el XOR acumulado de todos los bytes de la trama excepto el
propio byte de checksum.  El receptor verifica haciendo XOR de todos los
bytes incluyendo el checksum: si el resultado es \texttt{0x00}, la trama
es válida.

\section{Máquina de estados del decodificador}

El decodificador implementado en \class{UnerProtocol::feed()} procesa los
bytes de entrada uno a uno siguiendo los estados:

\begin{center}
WaitU $\to$ WaitN $\to$ WaitE $\to$ WaitR $\to$ ReadLen $\to$ WaitColon $\to$ ReadPayload
\end{center}

Al completarse \texttt{ReadPayload} con checksum correcto, se llama a
\texttt{processPayload()} que emite la señal \texttt{frameReceived}.

% ──────────────────────────────────────────────────────────────
\chapter{API de clases}
% ──────────────────────────────────────────────────────────────

\section{UnerProtocol}

Archivo: \file{comunicacion/UnerProtocol.h} / \file{.cpp}

\subsection{Encoder (métodos estáticos)}

\begin{tabular}{lp{8cm}}
  \toprule
  \textbf{Método} & \textbf{Descripción} \\
  \midrule
  \texttt{encode(cmd, payload)} & Framer genérico; genera la trama completa. \\
  \texttt{ackAlive()} & Genera ACK de heartbeat (0xF0). \\
  \texttt{cmdStart(s0,s1,s2)} & Genera trama CMD\_START (0x50). \\
  \texttt{cmdStop()} & Genera trama CMD\_STOP (0x51). \\
  \texttt{cmdReset()} & Genera trama CMD\_RESET (0x53). \\
  \texttt{cmdVelocidad(velIdx)} & Genera trama CMD\_VELOCIDAD (0x54). \\
  \texttt{cmdBlindDist(cfg,n)} & Genera trama CMD\_BLIND\_DIST (0x60). \\
  \texttt{cmdTrigger()} & Genera trama CMD\_TRIGGER (0x61). \\
  \texttt{cmdAnchoCaja(cm)} & Genera trama CMD\_ANCHO\_CAJA (0x62). \\
  \texttt{cmdCalibracion(cfg)} & Genera trama CMD\_CALIBRACION (0x63). \\
  \texttt{cmdHcsr04Cfg(cfg)} & Genera trama CMD\_HCSR04\_CFG (0x64). \\
  \texttt{cmdSg90Cfg(cfg)} & Genera trama CMD\_SG90\_CFG (0x65). \\
  \texttt{cmdIrDebounce(d)} & Genera trama CMD\_IR\_DEBOUNCE (0x66). \\
  \texttt{cmdTimersCfg(cfg)} & Genera trama CMD\_TIMERS\_CFG (0x67). \\
  \bottomrule
\end{tabular}

\subsection{Decoder}

\begin{tabular}{lp{8cm}}
  \toprule
  \textbf{Método / Señal} & \textbf{Descripción} \\
  \midrule
  \texttt{feed(data)} & Alimenta bytes al decodificador. \\
  \texttt{reset()} & Reinicia la máquina de estados. \\
  \texttt{frameReceived(f)} & Señal emitida al recibir una trama válida. \\
  \texttt{checksumError(raw)} & Señal emitida ante error de checksum. \\
  \bottomrule
\end{tabular}

\section{SerialManager}

Archivo: \file{comunicacion/SerialManager.h} / \file{.cpp}

\subsection{Métodos públicos}

\begin{tabular}{lp{7.5cm}}
  \toprule
  \textbf{Método} & \textbf{Descripción} \\
  \midrule
  \texttt{open(port, baud)} & Abre el puerto serie e inicia el heartbeat. \\
  \texttt{close()} & Cierra el puerto y detiene los timers. \\
  \texttt{isOpen()} & Devuelve si el puerto está abierto. \\
  \texttt{currentPort()} & Devuelve el nombre del puerto activo. \\
  \texttt{availablePorts()} & Lista los puertos COM disponibles. \\
  \texttt{sendStart(s0,s1,s2)} & Envía CMD\_START. \\
  \texttt{sendStop()} & Envía CMD\_STOP. \\
  \texttt{sendReset()} & Envía CMD\_RESET. \\
  \texttt{sendVelocidad(v)} & Envía CMD\_VELOCIDAD. \\
  \texttt{sendBlindDist(cfg,n)} & Envía CMD\_BLIND\_DIST. \\
  \texttt{sendTrigger()} & Envía CMD\_TRIGGER. \\
  \texttt{sendAnchoCaja(cm)} & Envía CMD\_ANCHO\_CAJA. \\
  \texttt{sendCalibracion(cfg)} & Envía CMD\_CALIBRACION. \\
  \texttt{sendHcsr04Cfg(cfg)} & Envía CMD\_HCSR04\_CFG. \\
  \texttt{sendSg90Cfg(cfg)} & Envía CMD\_SG90\_CFG. \\
  \texttt{sendIrDebounce(d)} & Envía CMD\_IR\_DEBOUNCE. \\
  \texttt{sendTimersCfg(cfg)} & Envía CMD\_TIMERS\_CFG. \\
  \bottomrule
\end{tabular}

\subsection{Señales}

\begin{tabular}{lp{7cm}}
  \toprule
  \textbf{Señal} & \textbf{Descripción} \\
  \midrule
  \texttt{connected(port)} & Puerto abierto exitosamente. \\
  \texttt{disconnected()} & Puerto cerrado. \\
  \texttt{connectionLost()} & Timeout de heartbeat. \\
  \texttt{errorOccurred(msg)} & Error en el puerto serie. \\
  \texttt{aliveReceived()} & Heartbeat 0xF0 recibido. \\
  \texttt{cajaMedida(cm)} & Caja medida (0x5F). \\
  \texttt{sensorIrActualizado(num,act)} & Estado sensor IR (0x5E). \\
  \texttt{brazoActuado(idx)} & Brazo actuado (0x52). \\
  \texttt{frameReady(f)} & Cualquier trama decodificada. \\
  \bottomrule
\end{tabular}

\section{Widgets}

\begin{tabular}{lp{9cm}}
  \toprule
  \textbf{Clase} & \textbf{Responsabilidad} \\
  \midrule
  \class{LedAliveWidget} & LED circular animado de heartbeat con estadísticas. \\
  \class{ConfigDialog} & Diálogo de calibración del clasificador (CMD 0x63). \\
  \class{VelocidadDialog} & Diálogo de ancho de caja de referencia (CMD 0x62). \\
  \class{AvanzadoDialog} & Diálogo de configuración avanzada de hardware (CMDs 0x64–0x67). \\
  \bottomrule
\end{tabular}

% ──────────────────────────────────────────────────────────────
\chapter{Implementación del protocolo -- detalles}
% ──────────────────────────────────────────────────────────────

\section{Corrección: máscara de bits en CMD\_BRAZO (0x52)}

El MCU envía en el byte \texttt{payload[0]} una \textit{máscara de bits}
donde el bit~$i$ indica que el servo~$i$ fue actuado:

\begin{lstlisting}[language=C++,caption={Despacho de CMD\_BRAZO en SerialManager}]
case Uner::CMD_BRAZO:
    if (frame.payload.size() >= 2) {
        const uint8_t mask = static_cast<uint8_t>(frame.payload.at(0));
        for (uint8_t i = 0; i < 3; i++) {
            if (mask & (1 << i))
                emit brazoActuado(i);
        }
    }
    break;
\end{lstlisting}

La versión anterior trataba \texttt{payload[0]} como un índice directo, lo
que causaba que el servo~2 (máscara = 4) se ignorara por superar el límite.

\section{CMD\_TRIGGER (0x61) -- trama sin payload}

El codificador llama a \texttt{encode(CMD\_TRIGGER)} sin payload.  Esto
genera \texttt{len = 2} (1 byte de cmd + 1 byte de checksum), que el MCU
interpreta correctamente como un comando sin datos.

\section{CMD\_BLIND\_DIST (0x60) -- envío parcial}

El parámetro \texttt{numBytes} (1–4) permite enviar solo el byte
\texttt{modo\_ciego} o también incluir las distancias \texttt{dist\_s0[0..2]},
facilitando la compatibilidad con versiones de firmware que no soportan
las distancias.

% ──────────────────────────────────────────────────────────────
\chapter{Estructura del repositorio}
% ──────────────────────────────────────────────────────────────

\begin{verbatim}
HMI_Clasificador/
  main.cpp                    Punto de entrada
  MainWindow.h / .cpp         Ventana principal
  HMI_Clasificador.pro        Archivo de proyecto qmake
  CLAUDE.md                   Guía de arquitectura para Claude Code
  Doxyfile                    Configuración de Doxygen
  comunicacion/
    UnerProtocol.h / .cpp     Protocolo UNER (encoder + decoder)
    SerialManager.h / .cpp    Gestor de puerto serie
  widgets/
    LedAliveWidget.h / .cpp   LED de heartbeat
    ConfigDialog.h / .cpp     Diálogo calibración (0x63)
    VelocidadDialog.h / .cpp  Diálogo ancho de caja (0x62)
    AvanzadoDialog.h / .cpp   Diálogo config. avanzada (0x64-0x67)
  docs/
    doxygen/                  Documentación generada por Doxygen
    informe/                  Este informe técnico
  Actividad-N-4-.../          Código fuente del firmware MCU (AVR C)
\end{verbatim}

% ══════════════════════════════════════════════════════════════
%  BIBLIOGRAFÍA / REFERENCIAS
% ══════════════════════════════════════════════════════════════
\chapter*{Referencias}
\addcontentsline{toc}{chapter}{Referencias}

\begin{itemize}
  \item Qt Documentation -- \href{https://doc.qt.io/qt-6/}{doc.qt.io/qt-6/}
  \item Qt Serial Port Module -- \href{https://doc.qt.io/qt-6/qtserialport-index.html}{QtSerialPort}
  \item Repositorio del proyecto -- \href{https://github.com/XradX314/HMI_Clasificador}{github.com/XradX314/HMI\_Clasificador}
  \item Documentación Doxygen -- \href{https://www.doxygen.nl/}{doxygen.nl}
\end{itemize}

\end{document}
